% Eine Aufgabe pro Textseite; die rechte Hälfte bleibt Arbeitsraum.
\Exercise{1}{Welche Mengen sind Schnitte?}
Alle folgenden Mengen bestehen aus rationalen Zahlen:
\[
\begin{aligned}
 A&=\varnothing,& B&=\mathbb Q,\\
 C&=\{p\in\mathbb Q:p\le1\},&
 D&=\{p\in\mathbb Q:p<1\},\\
 E&=\{p\in\mathbb Q:0<p<1\}.&&
\end{aligned}
\]
\textbf{(a)} Prüfe für jede Menge die vier Schnittbedingungen.
Lege eine Tabelle mit den Spalten \emph{nicht leer}, \emph{echt},
\emph{abwärts abgeschlossen} und \emph{ohne größtes Element} an.
Welche Mengen sind Schnitte?

\textbf{(b)} Gib für jede verletzte Bedingung einen konkreten Grund
oder ein Gegenbeispiel an. Warum ist es kein Widerspruch, dass
$\varnothing$ kein größtes Element besitzt?

\textbf{(c)} Zeichne $C,D,E$ auf rationalen Zahlengeraden.
Markiere besonders den Rand bei $1$. Welche Mehrdeutigkeit würde
entstehen, wenn man die vierte Schnittbedingung wegließe?
\ExerciseSpace{65}

\Exercise{2}{Eine rationale Zahl als Schnitt}
Sei $q\in\mathbb Q$ und
$\widehat q=\{p\in\mathbb Q:p<q\}$.

\textbf{(a)} Beweise alle vier Schnittbedingungen.
Finde für Nichtleerheit und Echtheit je eine konkrete rationale Zahl.
Ist $p<q$, so hilft für die letzte Bedingung die Zahl
\[
 r=\frac{p+q}{2}.
\]
Begründe ausdrücklich $p<r<q$.

\textbf{(b)} Welche der Zahlen $1,\frac32,2$ liegen in
$\widehat{\frac32}$? Finde zu $p=1$ ein größeres Element
von $\widehat{\frac32}$.

\textbf{(c)} Die Zahl $q$ ist eine obere Schranke von
$\widehat q$ \emph{innerhalb der rationalen Zahlen}.
Warum ist sie trotzdem kein größtes Element von $\widehat q$?
Erkläre den Unterschied zwischen \emph{größtes Element}
und \emph{oberer Schranke}.
\ExerciseSpace{70}

\Exercise{3}{Die Ordnung bleibt erhalten}
\textbf{(a)} Vergleiche
$\widehat{-\frac12}$ und $\widehat{\frac23}$ durch Inklusion.
Gib eine rationale Zahl an, die nur zum größeren Schnitt gehört.

\textbf{(b)} Beweise für beliebige $p,q\in\mathbb Q$:
\[
 p\le q\quad\Longleftrightarrow\quad
 \widehat p\subseteq\widehat q.
\]
Für die Rückrichtung: Angenommen $p>q$.
In welchem der beiden Schnitte liegt dann die Zahl $q$?

\textbf{(c)} Folgere:
\[
 p<q\quad\Longleftrightarrow\quad
 \widehat p\subsetneq\widehat q,
 \qquad
 \widehat p=\widehat q\quad\Longleftrightarrow\quad p=q.
\]
Warum darf die Abbildung $q\mapsto\widehat q$ eine
\emph{Einbettung} der rationalen Ordnung heißen?
\ExerciseSpace{65}

\Exercise{4}{Ein Zeuge ordnet zwei Schnitte}
Seien $A,B$ beliebige Schnitte. Wir wissen zunächst nicht,
ob sie Hauptschnitte sind.

\textbf{(a)} Angenommen, es gibt ein $a\in A\setminus B$.
Zeige für jedes $b\in B$, dass $b<a$ gilt.
Was wäre bei $b=a$, was bei $a<b$?

\textbf{(b)} Nutze die Abwärtsabgeschlossenheit von $A$ und
folgere $B\subsetneq A$.

\textbf{(c)} Beweise nun für alle Schnitte $A,B$:
\[
 A\subseteq B\quad\text{oder}\quad B\subseteq A.
\]
Welche Fallunterscheidung führt zu dem Zeugen aus Teil~(a)?
Warum gilt eine solche Vergleichbarkeit nicht für beliebige
Teilmengen von $\mathbb Q$? Gib ein einfaches Gegenbeispiel an.
\ExerciseSpace{80}

\Exercise{5}{Die größere Menge genügt}
\textbf{(a)} Setze
$A=\widehat{-\frac12}$ und $B=\widehat{\frac23}$.
Bestimme $A\cup B$ und zeichne die drei Mengen.
Welche ist die kleinste obere Schranke von $\{A,B\}$
in der Schnittordnung?

\textbf{(b)} Beweise für beliebige Schnitte $A,B$, dass
$A\cup B$ einer der beiden Schnitte ist.
Benutze Aufgabe~4.

\textbf{(c)} Erkläre direkt durch Inklusion:
$A\cup B$ enthält $A$ und $B$; jede Menge, die beide enthält,
enthält auch $A\cup B$.
Warum ist $A\cup B$ damit das Supremum von $\{A,B\}$?
Warum ist es hier sogar ein Maximum?

\textbf{Weiterdenken.} Was ändert sich, wenn eine Familie
unendlich viele Schnitte enthält? Die nächste Aufgabe zeigt,
dass ihr Supremum dann kein Familienglied sein muss.
\ExerciseSpace{80}

\Exercise{6}{Ein Supremum ohne größtes Familienglied}
Betrachte die Familie
\[
 \mathcal A=\{\widehat r:r\in\mathbb Q,\ r<1\}.
\]
Hier ist $r$ eine rationale Zahl und $\widehat r$ ein Schnitt.

\textbf{(a)} Zeige, dass $\mathcal A$ nicht leer und in der
Schnittordnung durch $\widehat1$ nach oben beschränkt ist.

\textbf{(b)} Beweise durch beide Inklusionen:
\[
 \bigcup\mathcal A=\widehat1.
\]
Für die schwierigere Inklusion: Ist $p<1$, wähle
$r=(p+1)/2$. In welchem Familienglied liegt $p$?

\textbf{(c)} Zeige ohne Vorgriff auf den allgemeinen
Vollständigkeitssatz, dass $\widehat1$ die kleinste obere Schranke
von $\mathcal A$ ist. Nutze dazu die Eigenschaft der Vereinigung.

\textbf{(d)} Warum gehört $\widehat1$ nicht zu $\mathcal A$?
Finde zu jedem Familienglied ein echt größeres Familienglied.
Erkläre damit: Ein Supremum kann existieren, obwohl ein Maximum fehlt.
\ExerciseSpace{60}

\Exercise{7}{Warum zwei Voraussetzungen nötig sind}
\textbf{(a)} Sei
\[
 \mathcal H=\{\widehat q:q\in\mathbb Q\}.
\]
Beweise $\bigcup\mathcal H=\mathbb Q$.
Finde dazu für jedes $p\in\mathbb Q$ ein passendes Familienglied.
Warum ist diese Vereinigung kein Schnitt?

\textbf{(b)} Beweise, dass $\mathcal H$ keine obere Schranke
\emph{unter den Schnitten} hat. Wenn ein Schnitt $U$ alle
Hauptschnitte enthielte: Welche rationalen Zahlen müsste er enthalten?

\textbf{(c)} Nun sei $\mathcal E=\varnothing$ die leere
\emph{Familie} von Schnitten. Was ist $\bigcup\mathcal E$?
Warum ist jede reelle Zahl eine obere Schranke von $\mathcal E$?

\textbf{(d)} Zeige, dass die leere Familie trotzdem kein
reelles Supremum besitzt: Wähle zu einem beliebigen Schnitt $U$
ein $u\in U$ und vergleiche $\widehat u$ mit $U$.

\textbf{Merksatz.} Formuliere die beiden Voraussetzungen,
die diese Gegenbeispiele im Vollständigkeitssatz erklären.
\ExerciseSpace{55}

\Exercise{8}{Den Vollständigkeitsbeweis zusammensetzen}
Sei $\mathcal A\ne\varnothing$ eine Familie von Schnitten.
Es gebe einen Schnitt $U$ mit $A\subseteq U$ für alle
$A\in\mathcal A$. Setze $S=\bigcup\mathcal A$.

\textbf{(a)} Zeige zuerst, dass $S$ ein Schnitt ist.
Ordne jeder Zeile den entscheidenden Grund zu und fülle die Lücken:
\begin{enumerate}[label=\arabic*.]
 \item $S\subseteq\mathbb Q$, denn jedes Familienglied \ldots
 \item $S\ne\varnothing$: Wähle zunächst \ldots, dann \ldots
 \item $S\ne\mathbb Q$, denn $S\subseteq\ldots$
 \item Ist $p\in S$ und $q\in\mathbb Q$ mit $q<p$, so gibt es \ldots
 \item Ist $p\in S$, so liefert dasselbe Familienglied \ldots
\end{enumerate}

\textbf{(b)} Zeige danach, dass $S$ die kleinste obere Schranke
von $\mathcal A$ ist. Trenne die beiden Aussagen:
\[
 A\subseteq S\quad(A\in\mathcal A),
\]
\[
 \bigl(A\subseteq V\text{ für alle }A\in\mathcal A\bigr)
 \quad\Longrightarrow\quad S\subseteq V.
\]
Hier sei $V$ ein beliebiger Schnitt.

\textbf{(c)} Markiere in deinem Beweis genau die Stellen,
an denen du $\mathcal A\ne\varnothing$ und die obere Schranke $U$
brauchst. Weshalb kommt der Beweis ohne Folgen und Grenzwerte aus?
\ExerciseSpace{45}
